\chapter{Requirements}
\label{ch:requirements}

This chapter tells you what your computer needs before you install ProPrep,
and which extra programs you need for some of its tools. Most people need
nothing beyond the ProPrep installer itself. Read the table in
Section~\ref{sec:requirements-what} to see whether your project needs more.

\section{Your computer}
\label{sec:requirements-computer}

\begin{tabularx}{\linewidth}{@{}lL@{}}
\toprule
\textbf{Computer} & \textbf{What works} \\
\midrule
Mac with an Apple chip (M1, M2, \ldots) & Yes. This manual was written on one
(macOS 27). \\
Mac with an Intel chip & Yes, with its own installer file. \\
Linux & Yes, on 64-bit Intel or AMD processors (x86-64). The installer has
been tested on Rocky Linux. \\
Windows 10 or 11 & Yes, through the \term{Windows Subsystem for Linux} (WSL),
a part of Windows that runs a small Linux system. Windows 10 needs version
2004 or later. \\
\bottomrule
\end{tabularx}

\medskip
ProPrep does not need a fast computer to prepare a structure. Simulations are
different: they use every processor core you give them, and a long simulation
is best run on a computing cluster or a computer with a suitable graphics card
(see Section~\ref{sec:requirements-what}). For scale: on our test laptop
(11 processor cores, no graphics card), a shortened test simulation of
ubiquitin in water (15,729 atoms) took about one minute. ProPrep describes its
0.42\,ns Workshop protocol as taking about 41 minutes with \texttt{sander.MPI}
on 4 cores (Chapter~\ref{ch:runningmd}).

\paragraph{Disk space.} Plan for about 8\,GB of free space for the
self-contained installer (Option~A in Chapter~\ref{ch:installation}), or
about 10\,GB in total on Windows, including WSL. The conda installation
(Option~B) used 4.8\,GB on our test Mac. Your projects need extra space:
simulation trajectories grow quickly.

\paragraph{The Terminal.} You start ProPrep by typing one command in the
Terminal, the program where you type commands (on Windows, the Ubuntu
window of WSL). You do not need to know any other commands. Chapter~\ref{ch:installation}
shows how to open it.

\paragraph{A web browser.} The recommended way to use ProPrep,
\texttt{proprep-web}, runs inside your normal web browser (Safari, Chrome,
Firefox or Edge). The 3D structure viewer also opens in the browser. If you
prefer, you can use ProPrep entirely in the Terminal instead
(Chapter~\ref{ch:starting}).

\paragraph{An internet connection.} You need the internet to:
\begin{itemize}
  \item download the installer (about 1 to 1.5\,GB), or to install with conda;
  \item download structures from the PDB, AlphaFold and AlphaFill;
  \item run a BLAST homology search, which uses the NCBI web service
        (Chapter~\ref{ch:compare});
  \item open \texttt{proprep-web}: the browser page loads its terminal
        display from the internet;
  \item download some reference files the first time a tool needs them (for
        example, the chemical component dictionary used by the PDB Filter).
\end{itemize}
Without the internet you can still work on structure files you already have.
ProPrep also checks once a day whether a newer release exists; with no
connection it skips the check silently.

\section{What you need for what}
\label{sec:requirements-what}

ProPrep's installer brings ProPrep, \term{AmberTools} (the free Amber
programs, such as tLEaP, sander, cpptraj and antechamber) and the helper
programs used for protonation, membranes and docking. The table lists
everything else.

\begin{longtable}{@{}>{\raggedright\arraybackslash}p{0.27\linewidth}>{\raggedright\arraybackslash}p{0.18\linewidth}>{\raggedright\arraybackslash}p{0.47\linewidth}@{}}
\toprule
\textbf{What} & \textbf{Who needs it} & \textbf{What it is for} \\
\midrule
\endhead
ProPrep with AmberTools & Everyone & Everything in this manual. Comes with
the installer. \\
A web browser & Most people & \texttt{proprep-web} and the 3D structure
viewer. \\
Internet connection & Most people & Downloading structures, BLAST, the
\texttt{proprep-web} page. \\
\midrule
MODELLER license key & Optional & Filling missing loops and residues,
rebuilding missing atoms, standard mutations and homology models. MODELLER
itself is installed for you; it only needs a free academic key
(Section~\ref{sec:installation-modeller}). \\
Gaussian 16 & Optional & Making new parameters for ligands, modified amino
acids and metal sites (\key{2c}), and running Gaussian ONIOM QM/MM inputs
(\key{5c}). Gaussian is commercial software; you buy and install it
yourself. \\
ORCA & Optional & Writing and running ORCA QM/MM inputs (\key{5c}). ProPrep
needs ORCA's \texttt{orca\_mm} program even to write the files. \\
\texttt{pmemd} and a graphics card & Optional & Faster simulations.
\texttt{pmemd} is part of Amber, which is distributed separately from
AmberTools. \texttt{pmemd.cuda} also needs an NVIDIA graphics card (GPU).
Without them, use \texttt{sander}, which is included. \\
CREST and xtb & Optional & One way of refining dihedral parameters in
\key{2c}. Not needed for the other routes. \\
A computing cluster with SLURM & Optional & Running long simulations. ProPrep
can write the job scripts (\key{5b}). \\
\bottomrule
\end{longtable}

Appendix~\ref{app:software} lists every menu tool and what happens when the
extra program is missing.

\begin{needsoftware}
Without a MODELLER key, ProPrep still starts and every tool opens. Only the
repair of missing loops, residues and atoms, standard mutations, and homology
models are switched off. ProPrep prints a reminder each time it starts
(Section~\ref{sec:installation-modeller}).
\end{needsoftware}

\section{What you do not need}
\label{sec:requirements-not}

\begin{itemize}
  \item You do not need to install Python, conda or AmberTools separately when
        you use the self-contained installer. It brings its own copies.
  \item You do not need BLAST installed on your computer. ProPrep sends the
        search to NCBI's web service.
  \item You do not need administrator rights on a Mac or on Linux. ProPrep is
        installed in a folder in your home folder. On Windows you need them
        once, to switch on WSL.
\end{itemize}

\begin{watchout}
If another Amber or AmberTools installation is already set up on your
computer, ProPrep may use it instead of its own copy, without telling you.
Section~\ref{sec:installation-amberhome} shows how to check this.
\end{watchout}
